\section{Zadanie 6}
\autor{Kacper Orszulak}
Alicja i Bob rozgrywają pomiędzy sobą $n$ pojedynków. Alicja jest lepsza i każdy pojedynek wygrywa niezależnie od pozostałych z prawdopodobieństwem 0.6. Zwycięzcą jest osoba, która wygra więcej pojedynków. Za pomocą nierówności Czernowa ogranicz prawdopodobieństwo tego, że zwycięzcą jest Bob.

\subsection*{Rozwiązanie}
Niech $X_i$ to indykator tego, że Alicja wygrała z Bobem w $i$-tym pojedynku. Dostajemy $\mathcal{P}(X_i = 1) = 0.6 = p$. Niech $X = \sum_{i \in [n]} X_i$ oznacza ilość wygranych przez Alicję pojedynków. Okazuje się, że $X$ ma rozkład dwumianowy $B(n, p)$, a zatem $\mathbb{E}[X] = np$. Liczymy z funkcji tworzącej:
\[
    M_{X_i}(t) = \mathbb{E}[e^{tX_i}] = pe^{1 \cdot t} + (1-p) e^{0 \cdot t} = 1 + p(e^t-1) \le \footnote{Korzystamy z tego, że $1+x \le e^x$.} e^{p(e^t-1)}
\]
\[
    M_X(t) = \prod_{i = 1}^nM_{X_i}(t) \le \prod_{i=1}^n e^{p(e^t-1)} = e^{np(e^t-1)}
\]
Bob wygrywa dla $X \le \frac{n}{2}$. Policzmy zatem:
\[
    \mathcal{P}(X \le \frac{n}{2}) \overset{\text{Nierówność Czernowa}}{\le} \frac{\mathbb{E}[e^{tX}]}{e^{t\frac{n}{2}}} \le \frac{e^{np(e^t-1)}}{e^{t\frac{n}{2}}} = e^{n(pe^t - p - t\frac{1}{2})}
\]
Teraz minimalizujemy wykładnik:
\[
    f(t) = pe^t - p - t\frac{1}{2} \implies f'(t) = pe^t - \frac{1}{2} \text{ oraz } f''(t) = pe^t
\]
Dla $f'(t) = 0$ dostajemy $t = ln\left(\frac{5}{6}\right)$ oraz $\forall_t \; f''(t) > 0$, a więc $t = ln\left(\frac{5}{6}\right)$ minimalizuje wykładnik.

Dostajemy zatem:
\[
    \mathcal{P}(X \le \frac{n}{2}) \le e^{n\left(-\frac{1}{10}-\frac{ln\left(\frac{5}{6}\right)}{2}\right)} \approx e^{-n \; \cdot \; 0.00884}
\]


\subsection*{Odpowiedź}
\[
\boxed{\mathcal{P}(X \le \frac{n}{2}) \le e^{-n \; \cdot \; 0.00884}}
\]
